\documentclass[11pt,a4paper]{article}

\usepackage[a4paper,margin=1in]{geometry}
\usepackage{fontspec}
\IfFontExistsTF{Times New Roman}{\setmainfont{Times New Roman}}{\setmainfont{TeX Gyre Termes}}
\IfFontExistsTF{Menlo}{\setmonofont{Menlo}}{\setmonofont{Latin Modern Mono}}
\usepackage{amsmath,amssymb}
\usepackage{booktabs}
\usepackage{longtable}
\usepackage{array}
\usepackage{enumitem}
\usepackage{fancyhdr}
\usepackage[round,authoryear]{natbib}
\usepackage[hidelinks]{hyperref}
\usepackage{xurl}
\usepackage{microtype}

\setlength{\parindent}{0pt}
\setlength{\parskip}{0.65em}
\setlist[itemize]{leftmargin=1.4em}
\newcommand{\file}[1]{\path{#1}}
\newcommand{\code}[1]{\path{#1}}
\newcolumntype{L}[1]{>{\raggedright\arraybackslash}p{#1}}

\pagestyle{fancy}
\fancyhf{}
\lhead{\small Application Strength Gate}
\rhead{\small Dataset Readiness Planning}
\cfoot{\thepage}
\renewcommand{\headrulewidth}{0.4pt}

\title{\textbf{Application Strength / Dataset Readiness Gate}\\
\large Decision Report Before Surrogate Dataset Planning}
\author{Codex-assisted downstream gate report}
\date{June 17, 2026}

\begin{document}
\maketitle

\begin{abstract}
This report decides whether the project has enough application-strength evidence to move from
small stress analytics into surrogate dataset planning. The decision is
\code{READY_FOR_DATASET_PLANNING_WITH_LIMITATIONS}. This is deliberately not approval for
dataset construction, neural-surrogate training, machine-learning label export, broad parameter
sweeps, or production risk claims. The supported next step is a written surrogate dataset design
plan that preserves solver, boundary, Greek, and diagnostic cautions.
\end{abstract}

\tableofcontents
\newpage

\section{Purpose of the Application-Strength Gate}

The solver-validation ladder ended with Ticket 12's conservative decision
\code{PASS_SOLVER_VALIDATION_WITH_LIMITATIONS}. Pilot 01 then showed that the validated
baseline American Crank--Nicolson / projected SOR solver can produce small diagnostic stress-map
figures with metadata. The quantitative risk-analytics step added numbers for boundary movement,
continuation premium, Gamma concentration, PSOR effort, runtime, and LCP diagnostics.

This gate answers a narrower downstream question: are the application messages now clear enough
to justify a surrogate dataset design plan? It does not create a dataset. It does not create
labels. It does not train a neural surrogate. It is a human-review decision point between
diagnostic analytics and any later data-generation work.

The methodology is consistent with the project planning documents and the student methodology
handouts: first validate the classical free-boundary benchmark, then run risk analytics, and
only later consider supervised surrogate data and models
\citep{black_scholes_1973,merton_1973,brennan_schwartz_1977,wilmott_howison_dewynne_1995}.

\section{Evidence Reviewed}

The gate uses four evidence layers:
\begin{itemize}
  \item Ticket 12 solver-validation synthesis.
  \item Downstream design gate policy.
  \item Pilot 01 stress-map diagnostics.
  \item Quantitative risk-analytics tables and figures.
\end{itemize}

\begin{longtable}{L{0.26\textwidth} L{0.34\textwidth} L{0.28\textwidth}}
\toprule
Evidence source & Main evidence & Gate reading \\
\midrule
\endhead
Ticket 12 synthesis & \code{PASS_SOLVER_VALIDATION_WITH_LIMITATIONS}. & Baseline solver is credible for human-reviewed downstream planning, not automatic data generation. \\
Downstream design gate & Baseline CN/PSOR selected; Rannacher kept as comparison only. & Future design must not silently switch solvers. \\
Pilot 01 & 11 pilot runs; all \code{ACCEPTABLE_FOR_PILOT_REVIEW}. & Small stress-map diagnostics are feasible with required metadata. \\
Quantitative analytics & 31 controlled analytics rows; all \code{ACCEPTABLE_FOR_ANALYTICS_REVIEW}. & Application messages are now quantitative enough for design planning. \\
Methodology handouts & Risk analytics should precede dataset and surrogate work. & The next permitted step is dataset planning, not construction. \\
\bottomrule
\end{longtable}

\section{Supported Application Messages}

The strongest application messages are directional and diagnostic. They are not universal
theorems over all parameter regimes, but they are coherent with the methodology handouts and the
project's American-option intuition.

\begin{longtable}{L{0.30\textwidth} L{0.33\textwidth} L{0.25\textwidth}}
\toprule
Supported message & Evidence & Interpretation \\
\midrule
\endhead
American put boundary decreases as volatility rises. & At \(\tau/T=1\), boundary moneyness moves \(0.7668 \rightarrow 0.5667 \rightarrow 0.4000\) as \(\sigma=0.20,0.40,0.60\). & Higher volatility increases waiting value for the put in the tested regime. \\
Dividend-call boundary decreases as dividend yield rises. & At \(\tau/T=1\), boundary moneyness moves \(1.8979 \rightarrow 1.2666 \rightarrow 1.2332 \rightarrow 1.1666\) as \(q=0.03,0.08,0.10,0.14\). & Higher dividend yield makes early exercise more attractive in the tested call regime. \\
Higher volatility increases numerical effort. & Put mean PSOR iterations rise from \(6.99\) to \(29.42\); call mean iterations rise from \(7.35\) to \(14.92\). & Stress regimes can be more expensive and need iteration metadata. \\
LCP diagnostics remain small in tested analytics cases. & Max obstacle violation is \(0\); max equation violation is about \(4.89\times10^{-8}\); max complementarity product is about \(4.73\times10^{-10}\). & The diagnostic analytics are internally coherent for human review. \\
Gamma remains fragile and mask-dependent. & Full-grid max \(|\Gamma|=30\); strict-mask values are lower but still diagnostic. & Greek outputs require masks and should not become production labels yet. \\
\bottomrule
\end{longtable}

\section{Cautionary Findings}

The application evidence is strong enough for planning, but not strong enough for data generation.
The caution is especially important for boundaries and Greeks. The boundary curves are extracted
from the continuation premium \(U-\Phi\) using a threshold. They are useful diagnostics, not exact
analytical free-boundary solutions. Gamma is a finite-difference second derivative, so it is
especially sensitive near payoff kinks, maturity rows, and extracted boundary bands.

\begin{longtable}{L{0.29\textwidth} L{0.38\textwidth} L{0.21\textwidth}}
\toprule
Caution & Why it matters & Planning implication \\
\midrule
\endhead
Threshold boundary extraction & Boundary location depends on grid, threshold, interpolation, and time row. & Store threshold and found/not-found metadata. \\
Gamma fragility & Full-grid Gamma is dominated by kink and maturity behavior. & Use masks; do not export unmasked Gamma labels. \\
Representative cases only & The analytics use controlled one-at-a-time sweeps, not full regime coverage. & Dataset design must plan held-out regimes carefully. \\
Local runtime numbers & Runtime seconds are machine-specific. & Use runtime as planning evidence, not production benchmark. \\
Baseline solver retained & Rannacher remains comparison-only after Ticket 10A. & Label source must record solver variant explicitly. \\
\bottomrule
\end{longtable}

\section{Dataset-Readiness Decision}

\begin{center}
\fbox{\parbox{0.86\textwidth}{
\centering
\textbf{Decision:} \code{READY_FOR_DATASET_PLANNING_WITH_LIMITATIONS}\\
\textbf{Allowed next step:} write a surrogate dataset design plan.\\
\textbf{Not allowed yet:} generate datasets, labels, neural models, or broad parameter sweeps.
}}
\end{center}

\begin{longtable}{L{0.27\textwidth} L{0.53\textwidth}}
\toprule
Decision field & Value \\
\midrule
\endhead
Gate decision & \code{READY_FOR_DATASET_PLANNING_WITH_LIMITATIONS}. \\
Basis & Solver validation passed with limitations; pilot and analytics evidence show clear diagnostic messages. \\
Allowed next step & Surrogate dataset design plan only. \\
Main limitation & Application evidence is representative and diagnostic, not broad label readiness. \\
Blocked until later & Dataset construction, neural training, production Greek labels, broad sweeps, final claims across all regimes. \\
\bottomrule
\end{longtable}

\section{Allowed Next Step}

The next step may be a written surrogate dataset design plan. That plan should define the future
data schema, parameter ranges, metadata requirements, split strategy, diagnostic gates, and human
review criteria. It should still avoid generating arrays or labels until a separate construction
gate approves the design.

The plan should explicitly state that all future labels must come from the validated baseline
CN/PSOR solver unless a later method gate approves an alternative. If any Rannacher-smoothed
results are included later, they must be separately labelled and must not replace the baseline by
default \citep{rannacher_1984}.

\section{Blocked Actions}

\begin{longtable}{L{0.34\textwidth} L{0.46\textwidth}}
\toprule
Blocked action & Reason \\
\midrule
\endhead
Actual dataset construction & This gate approves planning only; no arrays or training tables should be exported. \\
Neural surrogate training & The project has not yet approved data schema, label policy, or held-out regime design. \\
Training-label export & Boundary and Greek quantities remain diagnostic and metadata-dependent. \\
Production Greek labels & Gamma remains fragile near kink, maturity, and boundary regions. \\
Broad regime sweeps & Current evidence comes from controlled pilot and analytics sweeps. \\
Final claims across all regimes & The tested cases are representative, not universal convergence proof. \\
Exact free-boundary claims & Boundaries are threshold-based diagnostics. \\
\bottomrule
\end{longtable}

\section{Preliminary Dataset Design Rules}

The following rules are not a dataset schema yet. They are constraints that the future dataset
design plan must respect.

\begin{longtable}{L{0.25\textwidth} L{0.33\textwidth} L{0.29\textwidth}}
\toprule
Rule area & Required planning rule & Evidence basis \\
\midrule
\endhead
Normalization & Use normalized strike \(K=1\) for the first design. & Pilot 01 and quantitative analytics. \\
Option families & Include American puts and dividend-paying American calls. & Both families produced coherent application messages. \\
Parameters & Plan controlled ranges for \(r\), \(q\), \(\sigma\), and \(T\). & Planning report and methodology handouts. \\
Solver identity & Use baseline CN/PSOR as the default source. & Ticket 12 and Ticket 10A. \\
Splits & Use regime-level train/test splits, not only random point splits. & Handout warns that surrogate credibility depends on held-out regimes. \\
Core quantities & Preserve value, payoff, continuation premium, and parameter metadata. & Solver and pilot reports. \\
Boundary metadata & Preserve threshold, interpolation, and found/not-found status. & Ticket 09 and analytics evidence. \\
Greek masks & Preserve kink, maturity, boundary, and strict-mask flags. & Ticket 10 and analytics Gamma caution. \\
Diagnostics & Require PSOR convergence and LCP residuals for every solved regime. & Ticket 08 and quantitative analytics. \\
\bottomrule
\end{longtable}

The machine-readable policy companion is:
\begin{center}
\file{reports/04_downstream/dataset_readiness_policy.csv}
\end{center}
That file is a planning policy table. It is not a dataset.

\section{Human Review Checklist}

\begin{longtable}{L{0.48\textwidth} L{0.33\textwidth}}
\toprule
Review question & Required answer before dataset construction \\
\midrule
\endhead
Is baseline CN/PSOR accepted as the initial label source? & Yes, or document a new method gate. \\
Are future parameter ranges small enough for a first dataset? & Yes, with explicit ranges and rationale. \\
Are held-out regimes defined before generation? & Yes. \\
Are boundary thresholds and interpolation rules fixed? & Yes, with metadata fields. \\
Are Gamma masks specified? & Yes, especially near \(S=K\), maturity, and boundary bands. \\
Are PSOR and LCP diagnostics required for every generated surface? & Yes. \\
Are dataset files still blocked until the construction plan is approved? & Yes. \\
Are neural models still blocked until after dataset review? & Yes. \\
\bottomrule
\end{longtable}

\section{Limitations}

\begin{itemize}
  \item The application analytics are small and controlled; they are not broad stress-regime
  coverage.
  \item Boundary curves are useful but threshold-based.
  \item Gamma remains diagnostic, not production-quality.
  \item Runtime evidence is local to the current machine and grid size.
  \item The methodology handouts describe later dataset and surrogate stages, but this gate does
  not implement those stages.
  \item The decision allows design planning only; it does not approve generation.
\end{itemize}

\section{Recommended Next Step}

The recommended next step is a surrogate dataset design plan. That plan should define the future
schema, parameter ranges, split strategy, output labels, mask fields, diagnostic acceptance
criteria, and human-review checkpoint. It should also identify optional higher-grid confirmation
cases before any large generation run.

The project should keep the current gate language visible: the application evidence is strong
enough to plan the next stage, but not strong enough to construct datasets or train neural
surrogates yet.

\section{Reproducibility and Source Record}

\begin{longtable}{L{0.34\textwidth} L{0.50\textwidth}}
\toprule
Item & Record \\
\midrule
\endhead
Ticket 12 synthesis & \file{reports/03_solver/tickets/ticket_12_solver_validation_synthesis.tex}. \\
Downstream gate & \file{reports/04_downstream/downstream_design_gate.tex}. \\
Pilot 01 report & \file{reports/04_downstream/pilot_stress_maps_report.tex}. \\
Quantitative analytics report & \file{reports/04_downstream/quantitative_risk_analytics_report.tex}. \\
Pilot tables & \file{results/02_pilot_stress_maps/tables/}. \\
Analytics tables & \file{results/03_quantitative_risk_analytics/tables/}. \\
Policy CSV & \file{reports/04_downstream/dataset_readiness_policy.csv}. \\
Bibliography & \file{reports/03_solver/references.bib}. \\
\bottomrule
\end{longtable}

\bibliographystyle{plainnat}
\bibliography{reports/03_solver/references}

\end{document}
